% ---------------------------------------------------------------------------
% Tema Beamer "archdati" — accento unico teal/indaco, font Lato, pdflatex.
% Colorato ma sobrio; niente slide divisorie vuote; footer con barra di
% avanzamento (per regolarsi sui tempi durante l'orale).
% ---------------------------------------------------------------------------
\usepackage[T1]{fontenc}
\usepackage[utf8]{inputenc}
\usepackage[italian]{babel}
\usepackage{csquotes}
\usepackage{lato}
\usepackage{inconsolata}
\renewcommand{\familydefault}{\sfdefault}
\usepackage{tikz}
\usepackage{pgfplots}
\pgfplotsset{compat=1.18}
\usepackage{booktabs}
\usepackage{colortbl}
\usepackage{array}

% --- Palette -----------------------------------------------------------------
\definecolor{teal}{HTML}{0E7C86}      % accento primario
\definecolor{ink}{HTML}{22314F}       % inchiostro (testo/titoli)
\definecolor{sfondo}{HTML}{FAF9F6}    % off-white caldo
\definecolor{amber}{HTML}{E0A458}     % highlight riservato (dato chiave / vincitore)
\definecolor{grigio}{HTML}{6B7280}
\definecolor{grigiochiaro}{HTML}{DED9D3}
\definecolor{pgblu}{HTML}{336791}     % identita' Citus/PostgreSQL
\definecolor{mongoverde}{HTML}{00684A}% identita' MongoDB

% --- Base --------------------------------------------------------------------
\usetheme{default}
\setbeamertemplate{navigation symbols}{}
\setbeamertemplate{blocks}[rounded][shadow=false]
\setbeamercolor{background canvas}{bg=sfondo}
\setbeamercolor{normal text}{fg=ink}
\setbeamercolor{structure}{fg=teal}
\setbeamercolor{frametitle}{fg=ink}
\setbeamercolor{itemize item}{fg=teal}
\setbeamercolor{itemize subitem}{fg=teal}
\setbeamercolor{enumerate item}{fg=teal}
\setbeamerfont{frametitle}{series=\bfseries,size=\Large}
\setbeamerfont{title}{series=\bfseries,size=\huge}
\usebeamercolor[fg]{normal text}

% --- Titolo di frame con riga d'accento --------------------------------------
\setbeamertemplate{frametitle}{%
  \vspace*{0.7em}%
  {\usebeamerfont{frametitle}\usebeamercolor[fg]{frametitle}\insertframetitle}\\[0.18em]%
  {\color{teal}\rule{2.4em}{2.4pt}}%
  \vspace*{0.1em}%
}

% --- Footer: sezione + n/tot + barra di avanzamento --------------------------
\setbeamertemplate{footline}{%
  \leavevmode%
  \hbox to \paperwidth{\hspace{1.4em}%
    {\footnotesize\color{grigio}\insertsection}\hfill
    {\footnotesize\color{grigio}\insertframenumber\,/\,\inserttotalframenumber}%
    \hspace{1.4em}}%
  \vskip3pt%
  \ifnum\inserttotalframenumber>0
    {\color{teal}\rule{\dimexpr\paperwidth*\insertframenumber/\inserttotalframenumber\relax}{2.2pt}}%
    {\color{grigiochiaro}\rule{\dimexpr\paperwidth-\paperwidth*\insertframenumber/\inserttotalframenumber\relax}{2.2pt}}%
  \else
    {\color{grigiochiaro}\rule{\paperwidth}{2.2pt}}%
  \fi
}

% --- Blocchi ----------------------------------------------------------------
% Blocco per uno dei due sistemi (colore identitario in intestazione).
\newenvironment{sysblock}[2]{%
  \setbeamercolor{block title}{bg=#1,fg=white}%
  \setbeamercolor{block body}{bg=#1!7,fg=ink}%
  \begin{block}{#2}}{\end{block}}

% Blocco "risultato / verdetto" con accento ambra.
\newenvironment{keyblock}[1]{%
  \setbeamercolor{block title}{bg=amber,fg=ink}%
  \setbeamercolor{block body}{bg=amber!12,fg=ink}%
  \begin{block}{#1}}{\end{block}}

% Evidenziazione inline del dato chiave.
\newcommand{\dato}[1]{{\color{teal}\bfseries #1}}
\newcommand{\vince}[1]{{\color{teal}\bfseries #1}}

% --- Codice (listings) ------------------------------------------------------
\usepackage{listings}
\lstset{
  basicstyle=\ttfamily\footnotesize,
  keywordstyle=\color{teal}\bfseries,
  commentstyle=\color{grigio}\itshape,
  stringstyle=\color{amber!80!black},
  backgroundcolor=\color{ink!6},
  frame=leftline, rulecolor=\color{teal}, framesep=6pt, xleftmargin=8pt,
  breaklines=true, showstringspaces=false, columns=fullflexible,
  captionpos=b, aboveskip=4pt, belowskip=2pt,
}
